\begin{table}[t]
\centering
\begin{threeparttable}
\caption{Summary Statistics}
\label{tab:summary}
\begin{tabular}{lrrr}
\toprule
 & N & Mean & SD \\
\midrule
\multicolumn{4}{l}{\textit{Panel A: Medicare Spending (2019)}} \\
Medicare std.\ spending/capita (\$) & 3,134 & 10,193.37 & 1,658.41 \\
Inpatient std.\ spending/capita (\$) & 3,134 & 2,702.02 & 488.65 \\
\midrule
\multicolumn{4}{l}{\textit{Panel B: Hospital Market Structure}} \\
Hospitals (all types) & 3,134 & 1.42 & 2.62 \\
HHI (equal-share) & 3,134 & 8,492.02 & 2,753.8 \\
Competitive ($\geq$ 2 hospitals) & 3,134 & 0.24 & 0.43 \\
\midrule
\multicolumn{4}{l}{\textit{Panel C: County Characteristics}} \\
FFS beneficiaries & 3,134 & 10,549.46 & 24,124.18 \\
Average HCC risk score & 3,134 & 0.96 & 0.11 \\
Population & 3,123 & 104,858.4 & 334,157.5 \\
Poverty rate & 3,123 & 0.14 & 0.06 \\
Share 65+ & 3,123 & 0.2 & 0.05 \\
Median age & 3,123 & 41.6 & 5.39 \\
Dual-eligible share (\%) & 3,124 & 0.19 & 0.08 \\
Black share (\%) & 1,942 & 0.08 & 0.11 \\
\bottomrule
\end{tabular}
\begin{tablenotes}[flushleft]
\item \textit{Notes:} Unit of observation is county. Medicare spending from CMS Geographic Variation PUF (2019, pre-COVID). Hospital counts from CMS Hospital Compare (2024), including Acute Care, Critical Access, and VA hospitals. HHI computed as equal-share Herfindahl--Hirschman Index ($10{,}000/N$). County demographics from ACS 2022 5-year estimates.
\end{tablenotes}
\end{threeparttable}
\end{table}
